\clearpage
\noindent\textit{Title:}\\
\textbf{Polarimetric MIMO Arrays for Automotive Radar}\\[0.25cm]
\textit{Author:} Ing. Martin Šimák\\[0.25cm]
\textit{Department:} Department of Microelectronics\\[0.25cm]
\textit{Research group:} Microwave Sensing, Signals and Systems (MS3)\\[0.25cm]
\textit{Promotor:} Prof. DSc. Olexander Yarovyi\\[0.25cm]
\textit{External supervisor:} Dr. Gabriel Schnoering\\[0.25cm]
\textit{Committee member:} Dr. Marco Spirito\\[0.25cm]
\textit{Committee member:} Ing. Jan Puskely, Ph.D.\\[0.25cm]
\textit{Abstract:}\\
Modern \glspl{adas} and autonomous driving systems rely heavily on millimetre-wave \gls{fmcw} radar operating at \qtyrange{77}{81}{\giga\hertz} due to its operational robustness under adverse weather and lighting conditions.
However, contemporary commercial sensors rely almost exclusively on single-polarized scalar \gls{mimo} arrays.
While effective for 3D point-cloud generation, these scalar architectures cannot capture vector scattering dynamics, creating a major bottleneck in distinguishing \glspl{vru} -- such as pedestrians and cyclists -- from surrounding urban clutter.
Although radar polarimetry encodes the full complex scattering matrix to reveal physical target mechanisms, directly transferring polarimetric decomposition techniques from satellite remote sensing to dynamic automotive radar is severely hindered by grazing incidence geometry, ground-plane multipath, near-field wavefront curvature, and virtual aperture phase distortions.

To overcome these domain transfer hurdles, this research proposal outlines a top-down, algorithm-driven engineering framework for the investigation, formalization, and synthesis of polarimetry-first \gls{mimo} array architectures.
The project unites front-end antenna design, 2D array topology optimization, and Doppler-resolved polarimetric signal processing across five core objectives: deriving hardware specifications via algorithm perturbation analysis, constructing a multiscale hybrid simulation engine, synthesizing dual-polarized array topologies, designing high-purity W-band front-ends, and validating target decomposition models through empirical radar measurement campaigns.

Ultimately, this work aims to transition polarimetric automotive radar from an unmanaged scalar approximation to a mathematically rigorous, dual-polarized spatial-Doppler framework.
By enabling high angular resolution alongside high-purity scattering matrix estimation, the proposed architecture provides the spatial separability required to isolate dynamic \gls{vru} limb motion from adjacent clutter.
This research will deliver quantitative front-end specification bounds, validated simulation tools, novel array topologies, and experimental proof-of-concept hardware, establishing the foundation for next-generation polarimetric sensors in intelligent transportation systems.